\begin{proof}
    Zunächst definieren wir für $x \in H \setminus \{ 0 \}$ den Unterraum $G_x \coloneqq \cl \varphi(\Aa)x \leq H$, so gilt $\varphi(A) G_x \leq G_x$.

    (i) $\implies$ (ii): Wähle $y \neq 0$ mit $\varphi(\Aa) y \neq \{ 0 \}$. Dann ist $G_y \neq \{ 0 \}$ und weil $\varphi$ irreduzibel ist folgt $G_y = H$ und $y$ ist zyklisch. Insbesondere gibt es für $x \in H \setminus \{0\}$ stets ein $a \in \Aa$ mit $\varphi(a) x \neq 0$. Damit ist $\varphi(\Aa) x \neq \{0\}$, damit $G_x = H$ und damit ist $\varphi(\Aa) x$ dicht in $H$.

    (ii) $\implies$ (i): Sei $G \leq H$ abgeschlossen mit $\varphi(\Aa) G \subseteq G$. Ist nun $G = \{0\}$, so ist nichts zu zeigen. Sonst gibt es ein $x \in G \setminus \{0\}$, womit folgt
    \[
        H = \cl \varphi(\Aa) x \subseteq \varphi(\Aa) G \subseteq G,
    \]
    also $G = H$.

    (iii) $\implies$ (iv): Folgt sofort, da jedes $T \in L_b(H)$ mit $I_H$ kommutiert.

    (iv) $\implies$ (iii): Es gilt
    \[
        \varphi(\Aa)^{\updownarrow} = \varphi(\Aa)^{\updownarrow\updownarrow\updownarrow} = L_b(H)^{\updownarrow} = \spn \{ I_H \},
    \]
    wobei die letzte Gleichheit gilt da ein Operator der mit jedem Operator kommutiert bereits ein Vielfaches der Identität sein muss (vgl. Übung).

    (iii) $\implies$ (i): Sei $G \leq H$ abgeschlossen mit $\varphi(\Aa) G \subseteq G$. Sei nun $P_G \in L_b(H)$ die orthogonale Projektion auf $G$, dann gilt für alle $a \in \Aa$
    \[
        \varphi(a) P_G = P_G \varphi(a) P_G.
    \]
    Anwenden auf $a^*$ und Adjungieren liefert
    \[
        P_G \varphi(a) = P_G \varphi(a) P_G,
    \]
    also ist $P_G \in \varphi(\Aa)^{\updownarrow} = \spn \{ I_H \}$. Also folgt $P_G = 0$ oder $P_G = I_H$, und damit $G = \{0\}$ oder $G = H$.

    (i) $\implies$ (iii): Sei $T \in \varphi(\Aa)^\updownarrow$, so gilt auch $T^* \in \varphi(\Aa)^\updownarrow$. Damit gilt also auch $T_r, T_i \in \varphi(\Aa)^\updownarrow$. Es reicht nun zu zeigen, dass sowohl $T_r$ als auch $T_i$ ein Vielfaches der Identität ist. Mit dem Spektralsatz erhalten wir
    \[
        T_r = \int t \dd E_r(t),
    \]
    wobei $E_r$ das Spektralmaß auf $\sigma(T_r)$ bezeichnet. Insbesondere gilt $E_r(\Delta) \varphi(a) = \varphi(a) E_r(\Delta)$, womit $\varphi(a) E_r(\Delta) x \in E_r(\Delta) H$ und damit $\varphi(a) E_r(\Delta) H \leq E_r(\Delta) H$. Setzen wir nun $G \coloneqq E_r(\Delta) H$, so folgt $G = \{0\}$ oder $G = \{H\}$, insbesondere ist $E_r(\Delta)$ also entweder identisch Null oder die Identität. Damit folgt (vgl. Übung) $\sigma(T_r) = \{ \lambda \}$ für ein $\lambda \in \R$. Über das Spektralintegral folgt $T_r = \lambda I_H$. Analog für $T_i$, womit $T \in \spn \{ I_H \}$ folgt.

    (v) $\implies$ (iii): Es gilt (vgl. Übung)
    \[
        \varphi(\Aa)^\updownarrow = (\cl_{\mathcal{T}_s} \varphi(\Aa))^\updownarrow = L_b(H)^\updownarrow = \spn \{ I_H \}.
    \]

    (ii), (iv) $\implies$ (v): Betrachte die $\ast$-Unteralgebra $\varphi(\Aa) + \spn \{ I_H \} \leq L_b(H)$, welche (trivialerweise) die Identität enthält. Nach einem Ergebnis aus der VU folgt nun
    \[
        \cl_{\mathcal{T}_s} (\varphi(\Aa) + \spn \{ I_H \})
        =
        (\varphi(\Aa) + \spn \{ I_H \})^{\updownarrow \updownarrow}
        \geq
        \varphi(\Aa)^{\updownarrow\updownarrow}
        =
        L_b(H),
    \]
    also sogar Gleichheit. Die zu zeigende Aussage folgt nun, wenn wir $I_H \in \cl_{\mathcal{T}_s} \varphi(\Aa)$ zeigen können. Dazu nehmen wir $x \in H \setminus \{0\}$ und $b \in \Aa$, so gilt
    \[
        \lim_{q \in \Aa_+} \varphi(\varepsilon(q))(\varphi(b) x)
        =
        \lim_{q \in \Aa_+} \varphi(\varepsilon(q) b) x
        =
        \varphi \left( \lim_{q \in \Aa_+} \varepsilon(q) b \right) x
        =
        I_H(\varphi(b) x).
    \]
    Zu $y \in H$ und $\delta > 0$ gibt es nun $b_\delta \in \Aa$ mit $\Vert y - \varphi(b_\delta) x \Vert_H \leq \delta$. Nach dem oben Gezeigten gibt es ein $q_0 \in \Aa_+$ mit $\Vert \varphi(\varepsilon(q)) \varphi(b_\delta) x - I_H(\varphi(b_\delta) x) \Vert \leq \delta$ für $q \geq q_0$. Also folgt
    \begin{align*}
        \Vert \varphi(\varepsilon(q)) y - I_H (y) \Vert
        &\leq
        \Vert \varphi(\varepsilon(q)) (y - \varphi(b_\delta) x) \Vert
        + \Vert \varphi(\varepsilon(q)) \varphi(b_\delta) x - I_H(\varphi(b_\delta) x) \Vert
        \\ & \quad + \Vert I_H(\varphi(b_\delta) x - y) \Vert
        \leq 3 \delta.
    \end{align*}
    Also konvergiert $\varphi(\varepsilon(q))$ in der starken Operatortopologie für $q \in \Aa_+$ gegen $I_H$.
\end{proof}

\begin{remark}
    Sei $\varphi$ eine nicht irreduzible Darstellung, dann gibt es also ein abgeschlossenes $G \leq H$ mit $\varphi(\Aa) G \subseteq G$. Wir betrachten nun den Operator $\restr{\varphi(a)}{G} \in L_b(G)$. Man verifiziert, dass dann auch $\restr{\varphi}{G} : \Aa \to L_b(G)$ eine Darstellung ist.

    Insbesondere können wir $\varphi(a)$ in eine Block-Diagonalmatrix zerlegen, wobei auf der Diagonale die Einschränkungen von $\varphi(a)$ auf $G$ und $G^\perp$ stehen.
\end{remark}

\begin{example}
    Sei $\Aa$ eine $C^*$-Algebra und $m : \Aa \to \C$ ein multiplikatives lineares Funktional, welches mit $\ast$ verträglich ist\footnote{Wenn $\Aa$ eine Eins hat so ist dies automatisch.}.

    Für $a \in \Aa$ definieren wir nun $\varphi(a) \coloneqq m(a) I_H \in L_b(H)$, so ist das entstehende $\varphi : \Aa \to L_b(H)$ eine Darstellung. Diese Darstellung ist irreduzibel genau dann, wenn $\dim H = 1$, was beispielsweise aus dem obigen Satz folgt.
\end{example}

\begin{corollary}
    Sei $\Aa$ eine kommutative $C^*$-Algebra und $H$ ein nicht-trivialer Hilbertraum. Ist $\varphi : \Aa \to L_b(H)$ eine Darstellung, dann ist $\varphi$ irreduzibel genau dann, wenn $\dim H = 1$. In dem Fall gilt $\varphi = m(\cdot) I_H$ für ein multiplikatives, lineares Funktional auf $\Aa$.
\end{corollary}

\begin{proof}
    Die Unteralgebra $\varphi(\Aa) \leq L_b(H)$ ist kommutativ, also ist $\varphi(\Aa)^{\updownarrow\updownarrow}$ ebenfalls kommutativ. Nun ist nach dem vorigen Resultat $\varphi$ genau dann irreduzibel wenn $\varphi(\Aa)^{\updownarrow\updownarrow} = L_b(H)$. Es ist $L_b(H)$ jedoch genau dann kommutativ, wenn $\dim H = 1$.

    In diesem Fall können wir $m(a)$ als den Skalaranteil von $\varphi(a)$ definieren, was auch den letzten Teil der Aussage liefert.
\end{proof}